
\section{Related Work}

\vspace{3pt}
\textbf{Quantization of LLMs.}
Quantization reduces the size of LLMs and speedup inference. There are two primary quantization strategies:
(1) \textbf{Weight-only quantization}~\cite{gptq,lin2023awq,dettmers2023spqr,kim2024squeezellm} benefits edge devices where the workload is memory-bound, improving weight-loading speed. However, for cloud services with high user traffic and required batch processing, this method falls short as it does not accelerate computation in compute-bound scenarios.
(2) \textbf{Weight-activation quantization} accelerates computation in batch processing by quantizing both weights and activations~\cite{llmint8,wei2022outlier,xiao2023smoothquant}. OmniQuant~\cite{OmniQuant} and Atom~\cite{zhao2023atom} exploring more aggressive quantizations (\textbf{\texttt{W4A4, W4A8}}) and mixed precision to enhance model quality and efficiency, though these can impact model accuracy and reduce serving throughput. QuaRot~\cite{ashkboos2024quarot} further refines \textbf{\texttt{W4A4}} by rotating weights and activations at the cost of increased computational overhead due to additional transformations required during inference.



\vspace{3pt}
\textbf{LLM serving systems.}
Numerous systems have been proposed for efficient LLM deployment. Orca~\cite{orca} employs iteration-level scheduling and selective batching in distributed systems. vLLM~\cite{vllm} features virtual memory-inspired PagedAttention, optimizing KV cache management. SGLang~\cite{sglang} enhances LLM programming with advanced primitives and RadixAttention. LMDeploy~\cite{lmdeploy} offers persistent batching and blocked KV cache features to improve deployment efficiency. LightLLM~\cite{lightllm} manages GPU memory with token-wise KV cache control via Token Attention, increasing throughput. MLC-LLM~\cite{mlcllm} utilizes compiler acceleration for versatile LLM deployment across edge devices. TensorRT-LLM~\cite{trtllm} is the leading industry solution and the most important baseline in this paper.



\textbf{LLM Accelerators.} Transformers and LLMs have also generated considerable research interest in domain-specific accelerator design. Several works, such as $A^3$~\cite{ham20203}, ELSA~\cite{ham2021elsa}, and SpAtten~\cite{wang2021spatten}, have applied pruning techniques to the attention module, while GOBO~\cite{zadeh2020gobo} and EdgeBERT~\cite{tambe2021edgebert} have investigated quantization approaches. Additionally, DOTA~\cite{qu2022dota} introduces a lightweight, runtime detector for omitting weak attention connections, coupled with specialized accelerators for transformer inference. Apart from attention optimizations, STA~\cite{fang2022sta} leverages $N$:$M$ sparsity and specialized softmax module to reduce off-chip communication. Moreover, DFX~\cite{hong2022dfx} exploits model parallelism and optimized dataflow for low-latency generation. However, these accelerators have yet to be scaled up to recent LLMs with billions of parameters.
